\scanpage{46}
% archive sheet number: 45

\mspage{8}
Donc le diagramme squelettique

$(*)$\quad\diagram{\begin{tikzcd}[ampersand replacement=\&]
H^\bullet_{\underline D}(U,\xi)\arrow[dr] \& \& H^\bullet_{\underline D}(U',\xi)\arrow[dl,no head]\\
\& H^\bullet_{\underline D}(U\cap U',\xi) \&
\end{tikzcd}},

de type $\underline\Psi^{\mathrm{opp}}$, \inserted{provient-il} d'un diagramme concret

\[
\gamma_0\deleted{\unsure{\xi|\varphi}}\overset{\mathrm{déf}}{=}\bar\varphi_*(\unsure{a|\xi}_{\bar X})\in\underline D(\underline\Psi).
\]

Considérons l'inclusion

$\underline\Psi\underset{*}{\overset{\underline k}{\subset}}\Delta^1\times\Delta^1={}$\diagram{\begin{tikzcd}[ampersand replacement=\&]
\& (0,1)\arrow[dr,no head] \&\\
(0,0)\arrow[ur,no head]\arrow[dr,no head] \& \& (1,1)\\
\& (1,0)\arrow[ur] \&
\end{tikzcd}}

déduite de $\Psi$ en lui rajoutant l'élément final manquant $(1,1)$. Soit $\underline k$ l'inclusion, et considérons le ``carré cartésien''

\[
\gamma=\underline k_*(\gamma_0)\in\underline D(\Delta^1\times\Delta^1)
\]

\deleted{Donc les fibres de $\gamma$ en les points de} On a

\[
\underline k^*(\gamma)\simeq\gamma_0
\]

donc \unsure{le|la} \deleted{fibre de $\gamma$ en} squelett. induit par $\gamma$ sur $\underline\Psi$ n'est autre que $(*)$ ci-dessus. Mais je dis que le diagramme squelettique sur $\Delta^1\times\Delta^1$ tout entier s'identifie à

\mspage{9}
$(**)$\quad\diagram{\begin{tikzcd}[ampersand replacement=\&]
\& H^\bullet_{\underline D}(\deleted{\unsure{U}}\bar X,\xi)\arrow[dl,"i^*"' ]\arrow[dr,"i'{}^*"] \&\\
H^\bullet_{\underline D}(U,\xi)\arrow[dr,"i'_0{}^*"' ] \& \& H^{\unsure{\prime|\bullet}}_{\underline D}(U',\xi)\arrow[dl,"i_0^*"]\\
\& H^\bullet_{\underline D}(U\cap U',\xi) \&
\end{tikzcd}}

% slanted
\margin{\uline{NB} On a un énoncé analogue, pour \uline{tt} $\xi\in\underline D(X)$, pas seulement $\xi$ constant\ldots}

Notons pour ceci que si \deleted{que}

\[
\psi=\underline k\circ\unsure{y|\varphi}\colon\bar X\longrightarrow\Delta^1\times\Delta^1=\underline Q
\]

\deleted{\unsure{on}} $\gamma=\underline k_*(\gamma_0)=\underline k_*(\varphi_*(\xi_{\bar X}))=\psi_*(\xi_{\bar X})$, d'où

\[
H^\bullet_{\underline D}(\Delta^1\times\Delta^1,\gamma)\simeq H^\bullet_{\underline D}(\bar X,\xi)\xrightarrow{\sim}H^*_{\underline D}(X,\xi),
\]

\margin{car $\alpha\colon X\to\bar X$ est $\in W$}

D'autre part, comme $e=(1,1)$ élément final de $\Delta^1\times\Delta^1=\underline Q$, donc $e\to\underline Q$ \unsure{coasph.|asph.},

\[
\deleted{\unsure{H}}\quad\underbrace{H^\bullet_{\underline D}(Q,\gamma)}_{H^*_{\underline D}(\bar X,\xi)}\simeq\gamma(e)=\gamma(1,1)\deleted{\unsure{\in A}},
\]

d'où l'assertion.

Ça doit faire partie, d'autre part, des propriétés générales des carrés cartésiens \deleted{que} \inserted{``dans $A$''} (comme les), que

% slanted
\margin{\uuline{à vérifier}, car à mettre \unsure{en axiomes|en revenant|comme axiome} sur les dérivateurs.}

\[
i_0^*\ \text{iso}\Longleftrightarrow i^*\ \text{iso},
\]

Or comme c'est vrai pour \uline{tt} $\xi\in A,=\underline D(e)$, cela donne

\[
i_0\in W_{\underline D}\Longrightarrow{i}\in W_{\underline D}.
\]
